% --- CAPÍTOL 1: INTRODUCCIÓ ---
\chapter{Introducció i Motivació}
\thispagestyle{fancy}

\section{Motius de la tria}
La importància que té la comunicació per al nostre desenvolupament i la frustració per culpa de les dificultats que poden tenir algunes persones al intentar comunicar-se.

Comunicar-nos no només serveix per saber els interessos o conèixer a altres persones, sinó també per crear comunitats o grups socials. A més de tot això, l'ésser humà és un ser sociable i és frustrant quan la comunicació costa o hi ha dificultats.

\section{Orientació del treball}
Aquest treball s'orienta cap a quatre eixos principals:
\begin{itemize}
    \item \textbf{Inclusió:} Enfocada a l'autonomia i la qualitat de vida.
    \item \textbf{Logopèdia:} Des d'una perspectiva terapèutica.
    \item \textbf{Educació especial:} Aplicació en l'àmbit de les escoles.
    \item \textbf{Vida diària:} Utilitat pràctica i quotidiana.
\end{itemize}

\section{Altres consideracions}
Existeix una necessitat real de millorar els recursos per a persones amb dificultats comunicatives. Cal tenir en compte l'impacte social i emocional de no poder comunicar-se, així com les possibles mancances actuals en aquest camp.
